\documentclass[a4paper,12pt]{article}
\usepackage[T2A]{fontenc}
\usepackage[utf8]{inputenc}
\usepackage[russian]{babel}
\usepackage{cmap}
\usepackage{amsmath,amssymb,mathtools}
\usepackage{helvet}
\renewcommand{\familydefault}{\sfdefault}
\usepackage{icomma}
\usepackage[a4paper,margin=19mm,top=18mm,bottom=18mm]{geometry}
\usepackage{microtype}
\usepackage{xcolor}
\usepackage{tikz}
\usetikzlibrary{arrows.meta,positioning,shapes.geometric,calc}
\usepackage{tabularx,booktabs,longtable,array}
\usepackage{enumitem}
\usepackage{fancyhdr}
\usepackage{setspace}
\usepackage[hidelinks]{hyperref}
\usepackage{parskip}
\usepackage{needspace}

\definecolor{accent}{HTML}{247B7B}
\definecolor{darkaccent}{HTML}{155555}
\definecolor{textgray}{HTML}{30363A}
\definecolor{softgray}{HTML}{F2F5F5}
\definecolor{warn}{HTML}{8A5A20}

\setlength{\parindent}{0pt}
\setlength{\parskip}{0.55em}
\setlength{\headheight}{25pt}
\onehalfspacing
\setlist[itemize]{leftmargin=1.55em,itemsep=0.25em,topsep=0.25em}
\setlist[enumerate]{leftmargin=1.8em,itemsep=0.5em,topsep=0.4em}
\renewcommand{\arraystretch}{1.25}

\pagestyle{fancy}
\fancyhf{}
\fancyhead[L]{\small\color{textgray}Математическое ожидание}
\fancyhead[R]{\small\color{textgray}Ксения Клыкова}
\fancyfoot[C]{\small\color{textgray}\thepage}
\renewcommand{\headrulewidth}{0.3pt}
\renewcommand{\headrule}{\hbox to\headwidth{\color{accent}\leaders\hrule height \headrulewidth\hfill}}

\newcommand{\LessonDate}{28.09.26}
\newcommand{\E}{\mathsf{E}}
\newcommand{\Prb}{\mathsf{P}}
\newcommand{\sectionline}{\par\vspace{0.2em}{\color{accent}\hrule height 0.6pt}\vspace{0.5em}}
\newcommand{\Key}[1]{\textbf{\color{darkaccent}#1}}
\newcommand{\smallnote}[1]{\begin{quote}\small\color{textgray}#1\end{quote}}

\begin{document}

% --- title page ---
\thispagestyle{empty}
\begin{center}
\vspace*{2.5cm}
{\Large\color{darkaccent}\textbf{ЧЕК-ЛИСТ}}\\[0.8cm]
{\fontsize{28}{34}\selectfont\textbf{Математическое ожидание}}\\[0.45cm]
{\Large дискретная случайная величина и типовые задачи}\\[1.5cm]

{\large Ксения Клыкова}\\[0.4cm]
{\normalsize \LessonDate}\\[2.0cm]

\begin{minipage}{0.72\textwidth}
\centering
\color{textgray}
Главная цель: научиться распознавать, \emph{что именно усредняется}, правильно строить распределение и быстро выбирать специальную формулу там, где полная таблица не нужна.
\end{minipage}

\vfill
{\small\color{textgray}Лёвин Артём Александрович эксклюзивно для Ксении Клыковой}

\vspace{0.8cm}
\begin{tikzpicture}[remember picture,overlay]
  \draw[accent!65,line width=1.2pt]
    ($(current page.south west)+(18mm,15mm)$)
    .. controls ($(current page.south west)+(60mm,28mm)$) and ($(current page.south east)+(-70mm,4mm)$)
    .. ($(current page.south east)+(-18mm,17mm)$);
\end{tikzpicture}
\end{center}
\clearpage

\section*{1. Что нужно понимать про математическое ожидание}
\sectionline

Пусть дискретная случайная величина $X$ принимает значения $x_1,x_2,\ldots,x_n$ с вероятностями $p_1,p_2,\ldots,p_n$. Тогда
\[
\boxed{\E(X)=x_1p_1+x_2p_2+\dots+x_np_n=\sum_{i=1}^{n}x_ip_i.}
\]

\Key{Смысл.} Математическое ожидание --- это средний результат, к которому приближается среднее арифметическое большого числа независимых повторений одного и того же эксперимента при неизменных условиях.

\Key{Очень важно.} $\E(X)$ не обязано быть наиболее вероятным значением и вообще не обязано входить в список возможных значений $X$.

\smallnote{Например, если $X$ --- число автоматов, в которых к вечеру закончится кофе, то $X$ может принимать только $0$, $1$ или $2$. Значение $\E(X)=0{,}8$ вполне корректно: это долгосрочное среднее, а не результат одного конкретного дня.}

\subsection*{Перед вычислением: три быстрые проверки}
\begin{itemize}
  \item Все вероятности должны удовлетворять $0\le p_i\le 1$.
  \item Для полного распределения обязательно $p_1+\dots+p_n=1$.
  \item Каждому значению $x_i$ должна соответствовать именно его вероятность $p_i$.
\end{itemize}

\subsection*{Мини-пример}
Пусть $X$ принимает значения $0,1,2$ с вероятностями $0{,}32$, $0{,}56$, $0{,}12$. Тогда
\[
\E(X)=0\cdot0{,}32+1\cdot0{,}56+2\cdot0{,}12=0{,}8.
\]

\clearpage

\section*{2. Как строить распределение: пример с двумя автоматами}
\sectionline

Пусть вероятность того, что кофе закончится в первом автомате, равна $0{,}2$, во втором --- $0{,}6$. \textbf{Предполагаем независимость} событий. Случайная величина $X$ --- число автоматов, в которых кофе закончится к вечеру.

\begin{center}
\begin{tabularx}{0.88\linewidth}{>{\centering\arraybackslash}p{0.16\linewidth}X>{\centering\arraybackslash}p{0.23\linewidth}}
\toprule
$X$ & Событие & Вероятность \\
\midrule
$0$ & кофе не закончится ни в одном автомате & $(1-0{,}2)(1-0{,}6)=0{,}32$ \\
$1$ & кофе закончится ровно в одном автомате & $1-0{,}32-0{,}12=0{,}56$ \\
$2$ & кофе закончится в обоих автоматах & $0{,}2\cdot0{,}6=0{,}12$ \\
\bottomrule
\end{tabularx}
\end{center}

Отсюда
\[
\E(X)=0\cdot0{,}32+1\cdot0{,}56+2\cdot0{,}12=0{,}8.
\]

\subsection*{Быстрый способ для суммы индикаторов}
Если $X$ --- число произошедших событий среди нескольких испытаний, то удобно представить
\[
X=I_1+I_2+\dots+I_n,
\]
где $I_k=1$, если $k$-е событие произошло, и $I_k=0$ иначе. Тогда
\[
\boxed{\E(X)=p_1+p_2+\dots+p_n.}
\]
Для двух автоматов сразу получаем $\E(X)=0{,}2+0{,}6=0{,}8$.

\Key{Плюс этого приёма:} для вычисления ожидания суммы индикаторов независимость не нужна. Независимость понадобится, если Вы хотите строить полное совместное распределение переменной $X$ через произведения вероятностей.

\subsection*{Если известны количества исходов}
Если из $N$ равновозможных билетов ровно $n_i$ дают выигрыш $x_i$, то
\[
p_i=\frac{n_i}{N},\qquad \E(X)=\sum x_i\frac{n_i}{N}.
\]
Не забывайте учитывать исход $X=0$, если часть билетов ничего не выигрывает.

\clearpage

\begingroup
\setstretch{1.02}
\setlength{\parskip}{0.12em}
\section*{3. Три специальные формулы, которые экономят время}
\sectionline

\subsection*{A. Число успехов в $n$ одинаковых испытаниях}
Если проводится $n$ испытаний, вероятность успеха в каждом испытании постоянна и равна $p$, а $X$ --- число успехов, то
\[
\boxed{\E(X)=np.}
\]

\Key{Пример.} Симметричную монету бросают $6$ раз. Ожидаемое число орлов:
\[
\E(X)=6\cdot\frac12=3.
\]
Полное распределение по формуле Бернулли можно построить, но для одного только математического ожидания это лишняя работа.

\subsection*{B. Число испытаний до первого успеха}
Пусть независимые испытания повторяются при постоянной вероятности успеха $p>0$. Если $T$ --- номер испытания, на котором впервые получен успех, то
\[
\boxed{\E(T)=\frac{1}{p}.}
\]

\Key{Пример.} Вероятность попадания при одном выстреле $0{,}05$. Среднее число выстрелов до первого попадания:
\[
\E(T)=\frac{1}{0{,}05}=20.
\]

\subsection*{C. Число испытаний до $r$ успехов}
При тех же условиях, если $T_r$ --- число испытаний до получения $r$ успехов (успехи не обязаны идти подряд), то
\[
\boxed{\E(T_r)=\frac{r}{p}.}
\]

\Key{Примеры.}
\begin{itemize}
  \item Монету бросают до двух орлов: $\E(T_2)=2/(1/2)=4$.
  \item Кубик бросают до двух шестёрок: $\E(T_2)=2/(1/6)=12$.
\end{itemize}

\smallnote{Число $12$ во втором примере --- среднее число бросков в длинной серии повторений такого эксперимента. Это не означает, что две шестёрки «должны» выпасть ровно за 12 бросков.}

\subsection*{D. Монета до появления обеих сторон}
После первого броска одна сторона уже появилась. Чтобы дождаться противоположной стороны, при $p=1/2$ в среднем нужно ещё $1/p=2$ броска. Поэтому
\[
\boxed{\E(T)=1+2=3.}
\]
\endgroup

\clearpage

% --- flowchart-only page ---
\thispagestyle{fancy}
\begin{center}
{\Large\bfseries\color{darkaccent}Алгоритм: как выбрать способ вычисления $\E$}
\end{center}
\vspace{0.6cm}

\tikzset{
  startstop/.style={ellipse,draw=accent,very thick,align=center,minimum width=42mm,minimum height=11mm,fill=softgray},
  process/.style={rounded corners=2mm,draw=accent,thick,align=center,text width=86mm,minimum height=12mm},
  decision/.style={diamond,aspect=2.4,draw=darkaccent,thick,align=center,text width=48mm,inner sep=1.5pt},
  arrow/.style={-{Latex[length=3mm]},thick,draw=textgray}
}

\begin{center}
\begin{tikzpicture}[node distance=9mm]
  \node[startstop] (start) {Определите, что означает случайная величина};
  \node[decision,below=of start] (fixed) {$X$ считает успехи при фиксированном числе испытаний?};
  \node[process,below=of fixed] (np) {Да: одинаковая вероятность $p$ в $n$ испытаниях --- $\E(X)=np$.\\Разные вероятности отдельных событий --- $\E(X)=\sum p_i$.};
  \node[decision,below=of np] (wait) {$X$ --- число испытаний до первого или нескольких успехов?};
  \node[process,below=of wait] (geom) {Да: до первого успеха $\E(T)=1/p$;\\до $r$ успехов $\E(T_r)=r/p$.};
  \node[process,below=of geom] (general) {Если специальная формула не подходит: выпишите значения $x_i$ и вероятности $p_i$, проверьте $\sum p_i=1$, затем вычислите $\E(X)=\sum x_ip_i$.};
  \node[startstop,below=of general] (end) {Проверьте смысл и единицы результата};

  \draw[arrow] (start) -- (fixed);
  \draw[arrow] (fixed) -- node[right,font=\small]{да} (np);
  \draw[arrow] (np) -- (wait);
  \draw[arrow] (wait) -- node[right,font=\small]{да} (geom);
  \draw[arrow] (geom) -- (general);
  \draw[arrow] (fixed.east) -- ++(27mm,0) |- node[pos=0.24,right,font=\small]{нет} (wait.east);
  \draw[arrow] (wait.east) -- ++(27mm,0) |- node[pos=0.24,right,font=\small]{нет} (general.east);
  \draw[arrow] (general) -- (end);
\end{tikzpicture}
\end{center}
\clearpage

\section*{4. Типичные ошибки и самопроверка}
\sectionline

\begin{tabularx}{\linewidth}{>{\bfseries\color{warn}}p{0.33\linewidth}X}
\toprule
Ошибка & Как исправить \\
\midrule
«Матожидание --- самый вероятный исход» & Думайте о среднем результате большой серии повторений. Мода и математическое ожидание --- разные характеристики. \\

Вероятности не дают сумму $1$ & Сначала восстановите недостающую вероятность, затем считайте $\E(X)$. \\

Перепутаны значения и вероятности & Запишите таблицу парами $(x_i,p_i)$ и только потом перемножайте. \\

Забыли нулевой выигрыш & В распределении должен быть исход $X=0$ с вероятностью всех проигрышных билетов. \\

Используют $1/p$ без условий & Формула требует повторяющихся независимых испытаний с постоянной вероятностью успеха $p>0$. \\

Путают «до двух успехов» и «два успеха подряд» & $r/p$ относится к $r$ успехам без требования подряд. Для серии подряд нужна другая модель. \\

Считают полное биномиальное распределение ради одного ожидания & Если нужен только средний счёт успехов, используйте $np$. \\
\bottomrule
\end{tabularx}

\subsection*{Самопроверка перед ответом}
\begin{enumerate}
  \item Что именно является случайной величиной и в каких единицах она измеряется?
  \item Какие значения она может принимать?
  \item Вероятности корректны и дают сумму $1$?
  \item Есть ли специальная формула, которая сокращает расчёт?
  \item Имеет ли итог разумный смысл? Например, для числа успехов в $n$ испытаниях ожидание должно лежать между $0$ и $n$.
\end{enumerate}

\clearpage

\section*{5. Тренировочный блок --- Путь А}
\sectionline

\textbf{10 задач.} Решения здесь не приводятся. Сначала определите случайную величину и выберите наиболее короткий корректный метод.

\begin{enumerate}
  \item Случайная величина $X$ принимает значения $-1$, $0$, $2$ с вероятностями $0{,}2$, $0{,}5$, $0{,}3$. Найдите $\E(X)$.

  \item Случайная величина $X$ принимает значения $0$, $1$, $4$. Известно: $\Prb(X=0)=0{,}35$, $\Prb(X=1)=0{,}45$. Найдите $\Prb(X=4)$ и $\E(X)$.

  \item В лотерее выпущено $20\,000$ билетов. Из них $900$ билетов выигрывают 100 руб., $80$ билетов --- 1000 руб., $20$ билетов --- 5000 руб.; остальные билеты ничего не выигрывают. Найдите математическое ожидание выигрыша на один билет.

  \item Вероятность того, что к вечеру кофе закончится в первом автомате, равна $0{,}3$, во втором --- $0{,}4$. События независимы. Пусть $X$ --- число автоматов без кофе. Постройте распределение $X$ и найдите $\E(X)$.

  \item Симметричную монету бросают $14$ раз. Найдите математическое ожидание числа выпавших орлов.

  \item Игральный кубик бросают $30$ раз. Найдите математическое ожидание числа выпадений шестёрки.

  \item Вероятность успешного исхода в одном независимом испытании равна $0{,}2$. Испытания повторяют до первого успеха. Найдите ожидаемое число испытаний.

  \item Вероятность попадания в мишень при каждом выстреле равна $0{,}25$, выстрелы независимы. Стреляют до четырёх попаданий. Найдите ожидаемое число выстрелов.

  \item Симметричную монету бросают до тех пор, пока не появятся и орёл, и решка хотя бы по одному разу. Найдите ожидаемое число бросков.

  \item Три датчика срабатывают независимо. Вероятности срабатывания за смену равны соответственно $0{,}15$, $0{,}40$, $0{,}65$. Пусть $X$ --- число сработавших датчиков. Найдите $\E(X)$. Затем объясните одним предложением, почему для вычисления этого ожидания можно обойтись без полного распределения $X$.
\end{enumerate}

\clearpage

\section*{6. Ответы}
\sectionline

\begin{center}
\begin{tabularx}{0.92\linewidth}{>{\centering\arraybackslash}p{0.10\linewidth}X}
\toprule
№ & Ответ \\
\midrule
1 & $\E(X)=0{,}4$. \\
2 & $\Prb(X=4)=0{,}20$; $\E(X)=1{,}25$. \\
3 & $13{,}5$ руб. \\
4 & $\Prb(X=0)=0{,}42$, $\Prb(X=1)=0{,}46$, $\Prb(X=2)=0{,}12$; $\E(X)=0{,}7$. \\
5 & $7$. \\
6 & $5$. \\
7 & $5$. \\
8 & $16$. \\
9 & $3$. \\
10 & $1{,}20$; используется линейность математического ожидания: $\E(X)=0{,}15+0{,}40+0{,}65$. \\
\bottomrule
\end{tabularx}
\end{center}

\vfill
\small\color{textgray}
Ключевой ориентир: сначала определите случайную величину, затем проверьте, можно ли применить $np$, $1/p$, $r/p$ или сумму вероятностей индикаторов. Если нет --- стройте распределение и используйте $\E(X)=\sum x_ip_i$.

\end{document}
